% ===========================================================================
\section{Implementation}
\label{sec:impl}
% ===========================================================================

\system{} is implemented in approximately 3,200 lines of Python 3.12, organized
as a \texttt{uv} package with strict type checking (\texttt{ty check}) and
formatting enforcement (\texttt{ruff}).

\subsection{Fault Injection}
\label{sec:injection}

\system{} implements five fault injection mechanisms, all operating from
outside the container (no agent required inside):

\para{Network faults (\texttt{nw})} Injected using \texttt{nsenter} to enter
a container's network namespace, then applying \texttt{tc netem} rules:

\begin{lstlisting}[language=bash,caption={Network fault injection via nsenter.},label={lst:inject-nw}]
PID=$(docker inspect --format \
  '{{.State.Pid}}' $CONTAINER)
nsenter -t $PID -n tc qdisc replace \
  dev eth0 root netem delay ${DELAY}ms
\end{lstlisting}

This provides microsecond-precision delay injection without modifying container
images. Delay values are specified as continuous millisecond values (\eg{},
\texttt{slow-19.5ms}), enabling arbitrarily fine binary search granularity.

\para{Filesystem faults (\texttt{fs})} Simulated by injecting network delay on
the container's storage-path traffic. For higher fidelity, CharybdeFS
(FUSE-based) injection is supported but requires additional host setup.

\para{CPU faults (\texttt{cpu})} Docker's cgroup CPU quota
(\texttt{docker update --cpus=\$\{LIMIT\}}).
Severity values specify the fraction of a CPU core available.

\para{Memory faults (\texttt{mem})} Cgroup memory limits
(\texttt{docker update --memory=\$\{LIMIT\}}).
Values constrain available memory, forcing GC pressure in JVM-based systems.

\para{Process faults (\texttt{process})} Container restart or stop via
\texttt{docker restart}/\texttt{docker stop}.

\subsection{System Lifecycle}

Each trial executes a deterministic 8-step lifecycle: network creation,
container startup, stabilization wait (8--15\,s), fault injection, workload
execution, post-injection wait (12--30\,s), log collection, and teardown.
Each system is described by a YAML specification defining container image,
cluster size, start/stop commands, and workload command.

Table~\ref{tab:systems} lists all 10 systems with their cluster configurations.

\begin{table}[t]
\centering
\caption{Systems under test with cluster configurations.}
\label{tab:systems}
\small
\begin{tabular}{@{}llrl@{}}
\toprule
\textbf{System} & \textbf{Version} & \textbf{Nodes} & \textbf{Protocol} \\
\midrule
etcd & 3.5.10 & 3 & Raft \\
ZooKeeper & 3.8 & 3 & ZAB \\
MongoDB & 7.0 & 3 & Raft (replica set) \\
Redis & 7.2 & 6 & Gossip + PFAIL \\
TiKV & 7.5 & 3 & Multi-Raft \\
Cassandra & 4.1 & 3 & Gossip + $\phi$ detector \\
HBase & 2.5 & 4 & ZK-coordinated \\
Kafka & 3.6 & 3 & KRaft \\
CockroachDB & 23.2 & 3 & Multi-Raft \\
Hadoop & 3.3 & 3 & ZK-coordinated \\
\bottomrule
\end{tabular}
\end{table}

\subsection{Oracle Specifications}

We developed 21 oracle specifications covering multiple failure modes per
system: leader election (etcd, ZK, MongoDB, TiKV, CockroachDB), failover/PFAIL
(Redis), replication lag (Kafka, HBase), timeout cascades (Cassandra, HBase),
and throughput collapse (Hadoop, Cassandra).

\subsection{Severity Parsing}

Severity strings follow type-specific formats that the minimizer parses
into continuous numeric values for binary search:
Network (\texttt{slow-\{value\}\{unit\}} $\rightarrow$ ms),
CPU (\texttt{cpus-\{fraction\}} $\rightarrow$ cores),
Memory (\texttt{memory-\{value\}\{unit\}} $\rightarrow$ bytes).
After convergence, the numeric result is formatted back into canonical
strings for reproducibility.

\subsection{RunTrial Protocol}

The minimizer is decoupled from infrastructure through a \texttt{RunTrial}
protocol (Python \texttt{Protocol} class). This abstraction allows the same
minimizer to operate with \texttt{LiveRunner} (Docker + nsenter, used in
evaluation), \texttt{MockRunner} (deterministic, for unit tests), or future
\texttt{K8sRunner} (Kubernetes, for production clusters).
